% concatenated from l_infty_alg_conn.tex
\documentclass[12pt]{article}
\usepackage[margin=2in,includefoot,footskip=12pt,]{geometry}
\usepackage{amsfonts,amsmath, amsthm, amssymb,hyperref,fullpage}
\usepackage{graphicx}
\usepackage{listings}
\usepackage{xcolor,enumitem}
\usepackage{tikz}
\usepackage{algorithm}
\usepackage{tcolorbox}
\usepackage[noend]{algpseudocode}
\usetikzlibrary{arrows}
\usetikzlibrary{graphs,graphs.standard}
\usetikzlibrary{shapes.geometric}
\usetikzlibrary{svg.path}
\usepackage{bm}
\usetikzlibrary{patterns}

\newtheorem{observation}{Observation}[section]
\newtheorem{theorem}{Theorem}[section]
\newtheorem{lemma}{Lemma}[section]
\newtheorem{cor}{Corollary}[section]
\newtheorem{pro}{Proposition}[section]
\newtheorem{defn}{Definition}[section]
\newtheorem{remark}{Remark}[section]
\newtheorem{ex}{Example}[section]
\newtheorem{conj}{Conjecture}[section]
\DeclareMathOperator{\inte}{int}
\DeclareMathOperator{\diag}{diag}
\DeclareMathOperator{\spec}{spec}
\DeclareMathOperator{\sgn}{sgn}
\DeclareMathOperator{\diam}{diam}
\DeclareMathOperator{\tr}{Tr}
\DeclareMathOperator{\vol}{vol}
%\DeclareMathOperator{\Diag}{DIAG}

\newcommand{\affl}[3]{%
	\noindent #1,
	\textsc{#3}\\
	Email: \texttt{#2}\\[1.5pt]
}


\newcommand{\defeq}{\vcentcolon=}



\DeclareMathOperator{\ine}{In}
\newcommand*{\permcomb}[4][0mu]{{{}^{#3}\mkern#1#2_{#4}}}
\newcommand*{\perm}[1][-3mu]{\permcomb[#1]{P}}
\newcommand*{\comb}[1][-1mu]{\permcomb[#1]{C}}

\title{ On the $l_\infty$-analog of Algebraic Connectivity}

\author{ M. Rajesh Kannan,  Rahul Roy}

\date{\today}



\begin{document}
	\maketitle
	\begin{abstract} 
		
		The algebraic connectivity of a graph, defined as the second smallest eigenvalue of its Laplacian matrix, admits a well-known variational characterization involving the $\ell_2$-norm. Motivated by the recent introduction of its $\ell_\infty$-analogue by Andrade and Dahl, we investigate the graph parameter $\gamma(G)$, obtained by replacing the $\ell_2$-norm with the $\ell_\infty$-norm in the corresponding optimization problem. We establish a simple and explicit combinatorial formula expressing $\gamma(G)$ as the ratio of the order of the graph to its maximum transmission, thereby providing a direct graph-theoretic interpretation of the parameter. As a consequence, we obtain a polynomial-time algorithm based on breadth-first search, significantly simplifying the previously known linear programming approach. We prove that $\gamma(G)$ characterizes graph connectivity and completely characterize all $\ell_\infty$-Fiedler vectors as the vectors
		\[
		\left\{\pm\left(1-\gamma(G)d(u,\cdot)\right):u\in \mathcal{M}(G)\right\},
		\]
		where $\mathcal{M}(G)$ denotes the set of vertices of maximum transmission. Furthermore, we derive rms of several classical graph invariants, including the distance spectral radius, Wiener index, algebraic connectivity, and Cheeger constant. Finally, we establish a product formula for $\gamma(G)$ under Cartesian products of graphs, leading to explicit expressions for important graph families such as hypercubes, Hamming graphs, grid graphs, and torus graphs. 
		
		
		
		
	\end{abstract}
	
	{\bf AMS Subject Classification(2010):} 05C12, 05C50, 05C76.
	
	\textbf{Keywords.} Algebraic connectivity, Distance matrix, Transmission, Cartesian product, $l_\infty$-norm.
	\section{Introduction}
	The Laplacian matrices of graphs, along with their eigenvalues and eigenvectors, have found numerous applications in various domains, including combinatorial optimization, electrical networks, graph signal processing, data science, and machine learning. The second smallest eigenvalue of the Laplacian matrix provides a robust measure of the connectivity of the associated graph \cite{alg-con}.  It also serves as a lower or upper bound for several important graph parameters, such as the isoperimetric number, maximum cut, independence number, genus, and diameter \cite{mohar-lap-sur-1, mohar-eigen-comb-opti, trevisan_spec-partition}. The corresponding eigenvector is widely used in graph partitioning problems and data clustering, a method known as spectral clustering \cite{trevisan_spec-partition, ulrike-spectral}. In his seminal paper \cite{alg-con}, Fiedler coined the term algebraic connectivity for the second smallest eigenvalue of the Laplacian matrix. The algebraic connectivity of a graph $G$, denoted by $a(G)$,  has an interesting optimization characterization involving the $l_2$-norm in the objective function and the constraints. Recently, in \cite{enide-geir-2024}, Andrade and Dahl studied the counterparts of this optimization problem with respect to the $l_1$-norm and the $l_\infty$-norm. The corresponding optimal values are denoted by $b(G)$ and $\gamma (G)$, respectively. They showed that $b(G)$ has an interesting connection to the sparsest cut problem, implying that computing $b(G)$ is an NP-hard problem. They also proposed a linear programming formulation to compute $\gamma (G)$. 
	
	In this article, we provide an explicit combinatorial formula for $\gamma (G)$ in terms of the maximum row sum of the distance matrix of $G$. This leads to a polynomial-time combinatorial algorithm for computing $\gamma (G)$. Further, we study the connections between  $\gamma (G)$ and various graph parameters associated with $G$. In addition, we derive a formula for $\gamma(G)$ when $G$ is the Cartesian product of finitely many graphs.
	
	
	Let $G = (V(G), E(G))$ be an unweighted, undirected, and simple graph, where $V(G)$ is the set of vertices and $E(G)$ is the set of edges of $G$. Let the number of vertices in $G$, denoted by $\vert V(G) \vert $, be $n$. If the vertex $ u $ is adjacent to the vertex $v$, then we write $ u \sim v $.  The \textit{adjacency matrix} of a simple graph $ G $, denoted by $ A(G) $,  is the symmetric  $ n \times n $  matrix whose  $ (u,v)$-th entry is defined by $ a_{uv}=1 $ if $ u \sim v $, and  $ a_{uv} = 0 $ otherwise. The \textit{Laplacian matrix} $L(G)$ is the $n\times n$  matrix defined as 
	$$
	L(G)=\Delta(G)-A(G),
	$$
	where $\Delta(G)$ is the \textit{diagonal matrix} with the degrees of the vertices of $G$ as its diagonal elements, and $A(G)$ is the \textit{adjacency matrix} of $G$. The eigenvalues of adjacency matrices can be used to obtain various structural properties of graphs, such as bipartiteness, regularity, and more. Additionally, the eigenvalues of adjacency matrices provide bounds for several graph parameters, including the chromatic number, clique number, independence number, and others \cite{graphs-and-matrices, ratio-bound,  nica-book-2018, clique-indep-bound-wilf}.
	On the other hand, the eigenvalues of the Laplacian matrix of a graph reveal information about the connectedness of the graph $G$. In particular, $0$ is always an eigenvalue of $L(G)$, with the all-ones vector as a corresponding eigenvector. The second smallest eigenvalue of $L(G)$, denoted by $a(G)$, is called the \textit{algebraic connectivity} of the graph $G$ \cite{alg-con}. It is well known that $G$ is connected if and only if $a(G) > 0$ \cite{fiedler-1989}.
	An eigenvector corresponding to $a(G)$ is called a \textit{Fiedler vector} \cite{graphs-and-matrices}. The optimization approach to computing $a(G)$ is well known and is based on the Courant-Fischer theorem \cite{matrix-analysis}. A variational characterization of $a(G)$ is given as follows:
	\begin{equation*}
		a(G)=\min  \biggl\{\sum_{uv \in E(G)}(x_u - x_v)^2: \sum_{v \in V(G)} x_v =0 ~\mbox{and}~\Vert x \Vert _{2}=1\ \biggl\},
	\end{equation*} 
	where $\Vert x \Vert_2 := \sqrt{\sum\limits_{i=1}^n \vert x_i \vert^2}$ is the $l_2$-norm of the vector $x\in \mathbb{R}^n$. Note that the vector $x$ that attains the minimum in the above optimization problem has the following property: the entries of $x$ exhibit the minimum squared difference along the edges of the graph. This can be viewed as an "optimal" variation along the edges of $G$. For this reason, Andrade and Dahl \cite{enide-geir-2024} referred to this as the $l_2$-smoothing problem. Furthermore, they investigated analogous problems under the $l_1$-norm and the $l_\infty$-norm, considering both the objective function and the constraint \cite{enide-geir-2024}. This led to the introduction of two new parameters, defined as follows:
	$$
	b(G)=\min\biggl\{\sum_{uv \in E(G)}\vert x_u - x_v\vert: \sum_{v \in V(G)}x_v=0~\mbox{and}~ \Vert x \Vert_{1}=1\biggl\},
	$$
	$$
	\gamma(G)=\min\biggl\{\max_{uv \in E(G)}\vert x_u - x_v\vert: \sum_{v \in V(G)} x_v =0 ~\mbox{and}~\Vert x \Vert _{\infty}=1\biggl\}.
	$$
	Here, $\Vert x \Vert_1 := \sum\limits_{i=1}^n |x_i|$ denotes the $l_1$-norm, and $\Vert x \Vert_\infty := \max \{ |x_i| : 1 \leq i \leq n \}$ denotes the $l_\infty$-norm of the vector $x \in \mathbb{R}^n$. These optimization problems are well-defined, as the minima are attained due to the compactness of the constraint sets and the continuity of the objective functions. The vectors that give the optimal solutions are referred to as the {$l_1$-\textit{Fiedler vector} and the {$l_\infty$-\textit{Fiedler vector}, respectively. Furthermore, they formulated the following linear programming problem to compute  $\gamma(G)$. Let $G$ be a graph with $n$ vertices. For every $k$ with $1 \leq k \leq n$, consider the linear programming problem $LP(k)$:
			\begin{align*}
				\text{Minimize} \quad & y \\
				\text{subject to} \quad & x_i - x_j \leq y \quad  (ij \in E(G)), \\
				& -(x_i - x_j) \leq y \quad (ij \in E(G)), \\
				& \sum_{j=1}^{n} x_j = 0, \\
				& -1 \leq x_j \leq 1 \quad (1\leq j \leq n), \\
				& x_k = 1.
			\end{align*}
			Solving $LP(k)$ for each $1 \leq k \leq n$, and choosing the minimum value of $y$ obtained gives the value of $\gamma(G)$. 
			
			The primary objective of this paper is to provide an elegant combinatorial expression for the quantity $\gamma(G)$ in terms of the maximum row sum of the distance matrix of the graph $G$. Furthermore, in the case of trees, it suffices to consider only the maximum row sum corresponding to the pendant vertices. This is done in Section \ref{main-results}. Next, in Section \ref{compari-other}, we establish bounds for $\gamma(G)$ in terms of well-known graph-theoretic parameters, namely the distance spectral radius, Wiener index, algebraic connectivity, Cheeger constant, and others. We also study the corresponding extremal graphs. In Section \ref{sec-product}, we present the results on the Cartesian product of graphs in relation to $\gamma (G)$. 
			\section{Preliminaries}\label{sec_pre}
			In this section, we recall some of the known definitions and results that will be used later.
			For a vertex $u\in V(G)$, let $\deg(u)$ denote the degree of $u$ in $G$. A vertex $v$ of a graph $G$ is called a \textit{pendant vertex}, if $\deg(v) =1$. The distance between two vertices $u$ and $v$ is denoted by $d(u,v)$. The \textit{diameter} of a graph $G$, denoted by $\diam(G)$, is the maximum distance between any pair of vertices in $G$, that is,
			$$
			\diam(G)=\max_{ u,v \in V(G)} d(u,v).
			$$
			Let $u \in V(G)$. Define 
			$$
			S_r(u)=\biggl\{v \in V(G) : d(u,v)=r \biggl\}.
			$$
			The distance matrix $\mathcal{D}(G)$ of a graph $G$ is the $n \times n$ matrix in which the $(u,v)$-entry $\mathcal{D}_{uv}$ is given by the distance $d(u,v)$ between the vertices $u$ and $v$. The spectrum of $\mathcal{D}(G)$, denoted by  $\{\partial_1(G),\partial_2(G),\dots,\partial_n(G)\}$, consists of the eigenvalues of $\mathcal{D}(G)$ arranged in non-increasing order, that is, $\partial_1 (G) \geq \partial_2 (G) \geq \dots \geq \partial_n (G)$.
			
			An $n \times n$ matrix $A$ is said to be \textit{cogredient} to a matrix $E$ if there exists a permutation matrix $P$ such that $PAP^T=E$, where $P^T$ denotes the transpose of the matrix $P$. A matrix $A$ is called \textit{reducible} if it is cogredient to a block matrix of the form:
			$$
			E = \begin{bmatrix}
				B & 0 \\
				C & F
			\end{bmatrix},
			$$ where $B$ and $F$ are square matrices. Further, if $n=1$ and $A=0$, then $A$ is also considered reducible. Otherwise, $A$ is called \textit{irreducible}.
			The well-known Perron-Frobenius theorem states that if $A$ is non-negative and irreducible, then the spectral radius $\rho(A)$ is a simple eigenvalue of $A$   and any eigenvalue of $A$ that has the same modulus is also simple. Moreover, $A$ admits a positive eigenvector $x$ corresponding to $\rho(A)$, and every non-negative eigenvector of $A$ is a scalar multiple of $x$. For more details about the Perron-Frobenius theory, we refer to \cite{Non-negative-matrices}.  It is easy to see that the distance matrix $\mathcal{D}(G)$ of a connected graph $G$ is irreducible. Hence, by Perron-Frobenius theorem, we have $\partial_1(G)$ is simple and there is an eigenvector corresponding to $\partial_1(G)$ with strictly positive coordinates. 
			
			The \textit{transmission} of a vertex $u \in V(G)$, denoted by  $\tr(u)$, is the sum of distances from $u$ to all other vertices in $G$ \cite{Distance-spectra}. That is, 
			$$
			\tr(u)=\sum_{ v\in V(G)}d(u,v)
			$$
			A graph $G$ is said to be \textit{transmission-regular}, if $\tr(v)=k$  for some $k\in \mathbb{N}$ and for every vertex $v$ in $G$.
			A graph $G$ is called \textit{vertex-transitive} \cite{nica-book-2018} if, for any pair of vertices, there exists an automorphism of $G$ that maps one vertex to the other. In simpler terms, the structure of the graph looks the same from every vertex. These graphs are highly symmetric--- for example, the complete graph $K_n$, the cycle graph $C_n$, and the Petersen graph. 
			
			\section{A combinatorial formula for $\gamma (G)$}\label{main-results}
			In this section, we study the combinatorial significance of the quantity $\gamma(G)$. We begin by showing that $\gamma(G)$ characterizes the connectedness of a graph, as established in Theorem \ref{gamma-connected}.
			Then, in Theorem \ref{main theorem}, we prove that $\gamma(G)$ admits a neat combinatorial formula in terms of the transmission of vertices. Next, Theorem \ref{set of vectors} provides a complete characterisation of the $l_\infty$-Fiedler vectors.  We present a polynomial-time algorithm for computing $\gamma(G)$ based on the transmission values of the vertices of $G$. Furthermore, in Theorem \ref{pendant}, we show that for trees, it suffices to compute the transmission of the pendant vertices to determine $\gamma(G)$.
			The following notation will be used to derive the results presented in this section. Let $$\mathcal{F}=\biggl\{x \in \mathbb{R}^n:\sum_{v \in V(G)} x_v =0 ~\mbox{and} ~\Vert x \Vert_{\infty}=1\biggl\}.$$
			For a given $x \in \mathcal{F}$, define $$ \gamma_x(G)=\max_{uv \in E(G)}\vert x_u - x_v \vert .$$ In this notation, we have $$\gamma(G)=\min_{x \in \mathcal{F}}\gamma_x(G). $$ A vector  $x \in \mathcal{F}$ is called an $l_\infty$-\textit{Fiedler vector} of $G$ if $$\gamma(G)=\gamma_x(G).$$
			
			
			\begin{theorem}\label{gamma-connected}
				Let $G$ be a graph. Then $G$ is disconnected if and only if $\gamma(G)=0$.
			\end{theorem}
			\begin{proof}
				Let $G$ be a disconnected graph having two components $C_1$ and $C_2$. Without loss of generality,  let $|C_1|\leq |C_2|$. Consider the vector $x$
				
				$$   x_v =
				\left\{
				\begin{array}{ll}
					1  & \mbox{if } v \in V(C_1), \\
					-\frac{|C_1|}{|C_2|} & \mbox{if } v \in V(C_2).
				\end{array}
				\right.
				$$
				
				Then, it is easy to see that  $x \in \mathcal{F}$ and $\gamma_x(G) =0$. So $\gamma(G)=0$.
				
				Conversely, let $G$ be a graph with $\gamma(G)=0$. Let $x$ be an $l_{\infty}$-\textit{Fiedler vector}. Then $\gamma_x(G) =0$, which implies that $x_u=x_v$ whenever $uv \in E(G)$. 
				If $G$ is connected, then the vector $x$ should be a multiple of the all-ones vector.  However, this contradicts the assumption that $x\in \mathcal{F} $. Thus $\gamma (G) \neq 0.$
			\end{proof}
			Let $\mathcal{D}_M(G)$ denote the maximum row sum of $\mathcal{D}(G)$. Next, we prove the main result of this section.
			\begin{theorem}\label{main theorem}
				Let $G$ be a connected graph on $n$ vertices. Then
				$$
				\gamma(G)=\frac{n}{\max\limits_{u \in V(G)}  \tr(u)}=\frac{n}{\mathcal{D}_M(G)}.
				$$
				
				\begin{proof}
					Let $x \in \mathcal{F}$. Without loss of generality, let $x_u=1$ for some $u \in V(G)$. 
					Thus, for $v \in V(G)$, we have
					$$
					x_v \geq 1- \gamma_x(G) d(u,v).
					$$
					Consequently,
					$$
					0=\sum\limits_{v \in V(G)} x_v\geq \sum\limits_{v \in V(G)} (1- \gamma_x(G) d(u,v))=n- \gamma_x(G) \tr(u).
					$$
					Therefore
					$$
					\gamma_x(G) \geq \frac{n}{\tr(u)}\geq \frac{n}{\mathcal{D}_M(G)}.
					$$
					Next, let us prove the other side of the inequality holds for some vector. Let $u \in V(G)$ with $\tr(u)=\mathcal{D}_M(G)$. Define
					$$
					y_v=1 - \frac{n}{\mathcal{D}_M(G)}d(u,v).
					$$
					Then
					$$
					\sum\limits_{v\in V(G)}y_v=\sum\limits_{v\in V(G)}\left(1 - \frac{n}{\mathcal{D}_M(G)}d(u,v)\right)=n-\frac{n}{\mathcal{D}_M(G)}\tr(u)=0.
					$$
					Let $d=\diam(G)$. Choose a diametrical pair $a,b$. For every $v$,
					\begin{align*}
						d(a,v)+d(b,v) &\geq d\\
						\sum\limits_{v \in V(G)}(d(a,v)+d(b,v))&\geq nd\\
						\tr(a)+\tr(b) &\geq nd
					\end{align*}
					Also $\mathcal{D}_M(G) \geq \tr(a)$ and $\mathcal{D}_M(G) \geq \tr(b)$. Thus
					$$
					2\mathcal{D}_M(G) \geq nd.
					$$
					Therefore 
					\begin{equation}
						\mathcal{D}_M(G) \geq \frac{nd}{2}.\notag
					\end{equation}
					Since $d(u,v) \leq d$ for any $v\in V(G)$, we have,
					$$
					y_v=1 - \frac{n}{\mathcal{D}_M(G)}d(u,v)\geq 1-\frac{nd}{\mathcal{D}_M(G)}\geq 1-2=-1.
					$$
					Thus $y \in \mathcal{F}$.\\
					Finally, for an edge $vw$
					$$
					\vert d(u,v) -d(u,w) \vert \leq 1.
					$$
					\begin{align*}
						\vert y_v -y_w\vert &=\vert 1 - \frac{n}{\mathcal{D}_M(G)}d(u,v)-1+\frac{n}{\mathcal{D}_M(G)}d(u,w)\vert\\
						&\leq \frac{n}{\mathcal{D}_M(G)}\vert d(u,w) -d(u,v) \vert \\
						& \leq\frac{n}{\mathcal{D}_M(G)}.
					\end{align*}
					Therefore $ \gamma_y(G) \leq \frac{n}{\mathcal{D}_M(G)}$. Thus, $$\gamma(G)=\frac{n}{\max\limits_{u \in V(G)}  \tr(u)}=\frac{n}{\mathcal{D}_M(G)}.$$
				\end{proof}
			\end{theorem}
			Define
			$$
			\mathcal{M}(G)=\biggl\{u \in V(G): \tr(u)=\mathcal{D}_M(G)\biggl\}.
			$$
			In the following theorem, we fully characterise all $l_\infty$-Fiedler vectors associated with the graph $G$.
			\begin{theorem}\label{set of vectors}
				For every $u \in \mathcal{M}(G)$, define
				$$
				x^{(u)}_v=1-\gamma(G)d(u,v).
				$$
				Then the full set of $l_\infty$-Fiedler vectors is 
				$$
				\biggl\{\pm x^{(u)} : u \in \mathcal{M}(G) \biggr\},
				$$
				apart from repetitions when different maximum-transmission vertices produce the same vector.
				\begin{proof}
					Let $x$ be an $l_\infty$-Fiedler vector. Without loss of generality, let $x_u=1$ for some $u \in V(G)$. Then
					$$
					x_v \geq 1- \gamma(G) d(u,v).
					$$
					Using Theorem \ref{main theorem} and summing over all vertices, we obtain
					$$
					0\geq n-\gamma(G) \tr(u)=n(1-\frac{\tr(u)}{\mathcal{D}_M(G)})\geq 0.
					$$
					Therefore, equality holds everywhere. Hence 
					$$
					\tr(u)=\mathcal{D}_M(G)
					$$
					and 
					$$
					x_v = 1- \gamma(G) d(u,v)
					$$
					for every vertex $v \in V(G)$.
				\end{proof}
			\end{theorem}
			As a consequence of Theorem \ref{main theorem}, we present the following algorithm (in pseudocode) to compute the quantity $\gamma(G)$ for any arbitrary connected graph $G$.
			\begin{algorithm}[H]
				\caption{Algorithm to compute $\gamma(G) $}
				\textbf{Input:} A connected graph $G$ on $n$ vertices
				
				
				\textbf{Output:} $\gamma(G)$
				
				\begin{algorithmic}[1]
					\State $maxTransmission := 0$
					\For{each vertex $v_i$ in $V(G)$}
					\For{each vertex $v_j$ in $V(G)$}
					\State Compute $d(v_i, v_j)$ %\Comment{Shortest distance from $v_i$ to $v_j$}
					\EndFor
					\State $\tr(v_i) := \sum\limits_{j \ne i} d(v_i, v_j)$ %\Comment{Transmission of $v_i$}
					\If{$\tr(v_i) > maxTransmission$}
					\State $maxTransmission := \tr(v_i)$
					
					\EndIf
					\EndFor
					\State \Return $(\frac{n} {maxTransmission})$
				\end{algorithmic}
			\end{algorithm}
			\begin{remark} Since the distance between any two vertices in a connected graph can be computed in polynomial time, it follows that, by using Theorem \ref{main theorem}, $\gamma(G)$ can also be computed in polynomial time using the BFS technique.
			\end{remark}
			
			We next apply the formula given in Theorem \ref{main theorem} to determine $\gamma(G)$ for some well-known graph families.
			\begin{cor}\label{basic examples}
				\begin{enumerate}
					\item [(a)] Let $K_n$ denote the complete graph on $n$ vertices. Then, $\gamma(K_n)=\frac{n}{n-1}$.
					\item [(b)] Let $C_n$ denote the cycle graph on $n$ vertices. Then,
					$\gamma(C_n)=\frac{n}{\lfloor \frac{n^2}{4}\rfloor}$.
					\item[(c)] Let $S_n$ denote the star graph on $n$ vertices. Then $\gamma(S_n)=\frac{n}{2n-3}.$\label{star graph}
					
					
					
					\item[(d)]   Let $K_{m,n}$ denote the complete bipartite graph where the partite sets have sizes $m$ and $n$ with $m \geq n$. Then $\gamma(K_{m,n})=\frac{m+n}{2m+n-2}$.
					
					\begin{proof}
						\textbf{$(a)$} Since every vertex in $K_n$ is adjacent to every other vertex, $\tr(u)=n-1$ for every $u \in V(K_n)$, we get $\gamma (K_n) = \frac{n}{n-1}$.
						
						\textbf{$(b)$} Let $n=2l$ for some $l \in \mathbb{N}$ and $l \geq 2$. Let $u \in V(C_{2l})$ and $\vert S_r(u) \vert=a_r$. Then,
						$\diam(C_{2l})=l$ and $a_1=a_2=\dots=a_{l-1}=2$, $a_l=1$. Thus,
						$$
						\tr(u)=2\sum^{l-1}_{r=1}r+l=l^2.
						$$
						Hence $\gamma(C_{2l})=\frac{2l}{l^2}=\frac{2}{l}$.\\
						
						Let $n=2l+1$ for some $l \in \mathbb{N}$ and $l \geq 1$. Let $u \in  V(C_{2l+1})$. Similarly we have
						%     $\diam(C_{2l+1})=l$. Here $a_1=a_2=\dots=a_{l}=2$ and 
						$$
						\tr(u)=2\sum^{l}_{r=1}r=l(l+1),
						$$
						for every $u\in V(C_{2l+1})$. Hence $\gamma(C_{2l+1})=\frac{2l+1}{l(l+1)}$.
						
						Therefore, combining the results for even and odd values of $n$, we obtain $\gamma(C_n)=\frac{n}{\lfloor \frac{n^2}{4}\rfloor}$.
						\textbf{$(c)$} Let $u$ be the center vertex of $S_n$. Then $\tr(u)=n-1$. If $v$ is a pendant vertex, then $\tr(v)=2n-3$. So, $\mathcal{D}_M(S_n)=2n-3$. Thus $\gamma(S_n)=\frac{n}{2n-3}.$
						
						
						\textbf{$(d)$} Let $(V_1,V_2)$ be the bipartition of $K_{m,n}$ with $\vert V_1\vert=m$ and $\vert V_2\vert=n$. Then, for any vertex $u\in V_1$, $\tr(u)=n+2m-2$. Similarly, for any vertex $v \in V_2$, $\tr(v)=m+2n-2$. since $m\geq n$, $\mathcal{D}_M(K_{m,n})=2m+n-2$. Hence the result follows.
					\end{proof}
					
				\end{enumerate}
			\end{cor}
			For trees, we make the interesting observation that the maximum transmission occurs only at the pendant vertices. Therefore, to compute $\gamma(T)$ for a tree $T$, it suffices to consider only the transmissions at the pendant vertices.
			\begin{theorem} \label{pendant}
				Let $T$ be a tree, and let $u \in V(T)$ be such that $\tr(u)=\mathcal{D}_M(T)$. Then $u$ is a pendant vertex.
			\end{theorem}
			\begin{proof}
				Let $u \in V(T)$ be such that $\tr(u)=\mathcal{D}_M(T)$. Suppose that the vertex $u$ is not a pendant vertex. Let $\deg (u) \geq 2 $, and let $v$ and $w$ be two neighbours of $u$. Let $S_v$ and $S_w$ be the two branches rooted at $u$ containing the vertices $v$ and $w$, respectively. Without loss of generality,  assume that $\vert S_v \vert \geq \vert S_w \vert$.
				
				Let $p \in V(T)$ and $q \in V(T)$ belong to the branches $S_v$ and $S_w$, respectively. Then $d(w,p)=d(u,p)+1$ and $d(w,q)=d(u,q)-1$. Thus, 
				$$
				\tr(w)-\tr(u)\geq\vert S_v \vert - \vert S_w \vert+1 >0
				$$
				which is not possible, since $\tr(u)$ is maximum. Thus, $u$ must be a pendant vertex.
			\end{proof}
			Using Theorems \ref{main theorem} and \ref{pendant}, one can directly obtain the value of $\gamma(P_n)$, which was detailed in \cite{enide-geir-2024}.
			\begin{cor}
				Let $P_n$ denote the path graph on $n$ vertices. Then $\gamma(P_n)=\frac{2}{n-1}$.
			\end{cor}
			\begin{proof}
				Every pendant vertex of $P_n$ has transmission $\frac{n(n-1)}{2}$. Thus $\gamma(P_n)=\frac{2}{n-1}$. 
			\end{proof}
			\section{Comparison with other parameters}\label{compari-other}
			In this section, we compare the parameter $\gamma (G)$ with other well-known graph invariants, such as the spectral radius of the distance matrix, the Wiener index, algebraic connectivity,  $b(G)$, and the Cheeger's constant.
			
			For a connected graph $G$, let $\partial_1 (G)$ denote the spectral radius of the associated distance matrix $\mathcal{D}(G)$ of $G$.
			\begin{theorem}\label{max eigenvalue}Let $G$ be a connected graph on $n$ vertices. Then
				$$
				\gamma(G) \leq \frac{n}{\partial_1 (G)},
				$$
				where the equality holds if and only if  $G$ is a transmission-regular graph.
			\end{theorem}
			\begin{proof}
				Since the largest eigenvalue of a nonnegative matrix is bounded above by the maximum row sum \cite{matrix-analysis}, the inequality follows from Theorem \ref{main theorem}. 
				
				Suppose $G$ is a transmission-regular graph. Then for every vertex $u$ in $G$, we have $\tr(u)=k$ for some $k\in \mathbb{N}$. This implies that all row sums of the distance matrix $\mathcal{D}(G)$ are equal to $k$, so the all-ones vector is an eigenvector of $\mathcal{D}(G)$ corresponding to the eigenvalue $k$. Since the largest eigenvalue is bounded above by $k$, $k$ should be the largest eigenvalue. Thus $k=\mathcal{D}_M(G)= \partial_1 (G)$.
				
				Conversely, assume that $\mathcal{D}_M(G)= \partial_1 (G)$. By Perron-Frobenius theorem, let $z = (z_i)$ be a positive eigenvector corresponding to the eigenvalue $\partial_1 (G)$. Let $z_k=\max\limits_{i \in V(G)} z_i$. Then for the $k$-th row in $\mathcal{D}(G)$, we have
				$$
				\partial_1(G) z_k=\sum\limits_{j=1}^n D_{kj}z_j \leq z_k(\sum\limits_{j=1}^n D_{kj}).
				$$
				Hence $\partial_1(G) \leq \sum\limits_{j=1}^n D_{kj} $. Since $\mathcal{D}_M(G)= \partial_1 (G)$, we  have $\partial_1 (G) = \sum\limits_{j=1}^n D_{kj}$. Thus $z_j=z_k$ for all $j \in \{1,2, \dots,n\}$. Hence, the all-ones vector is an eigenvector corresponding to the eigenvalue $\partial_1$, that is, all the row sums of $\mathcal{D}_M(G)$ are equal. So, $G$ is a transmission-regular graph.
			\end{proof}
			\begin{remark}
				A number of lower bounds for $\partial_1(G)$ have been extensively studied in the literature \cite{Distance-spectra}.  These findings can be used to derive further upper bounds for $\gamma(G)$ across various classes of graphs.
			\end{remark}
			
			
			The \textit{Wiener index} $W(G)$ of a graph $G$ is the sum of distances between all unordered pairs of vertices of $G$ \cite{Distance-spectra}. Note that the Wiener index $W(G)$ of a graph $G$ is half of the sum of transmissions of all the vertices of $G$.
			$$
			W(G)=\frac{\sum\limits_{u \in V(G)}\tr(u)}{2}.
			$$
			
			In the next theorem, we establish a relationship between $\gamma (G)$ and the Wiener index of the graph $G$. 
			\begin{theorem}\label{weiner-index-bound}
				Let $G$ be a connected graph on $n$ vertices. Then
				$$
				\gamma(G) \leq \frac{n^2}{2W(G)}.
				$$
				Furthermore, equality holds if and only if  $G$ is a transmission-regular graph.
			\end{theorem}
			\begin{proof} 
				Note that
				\begin{align*}
					W(G)&=\frac{\sum\limits_{u \in V(G)}\tr(u)}{2} \\ 
					&\leq  \frac{n}{2} \max\limits_{u \in V(G)}\tr(u).
				\end{align*}
				
				That is,  $\max\limits_{u \in V(G)}\tr(u) \geq \frac{2W(G)}{n}$. The result now follows directly from Theorem \ref{main theorem}.
				
				The equality holds if and only if all vertex transmissions are equal; that is, $G$ is a transmission-regular graph.
			\end{proof}
			
			Next, we provide a bound for $b(G)$ in terms of $\gamma (G)$.  
			\begin{theorem}
				Let $G$ be a graph with $m$ edges. Then $b(G) \leq \frac{m}{2}\gamma(G)$.
			\end{theorem}
			\begin{proof}
				Let $x$ be an $l_{\infty}$-Fiedler vector of $G$. Let $u \in V(G)$ be such that $x_u=1$. Then
				$$
				\Vert x \Vert_1=1+\sum_{\substack{v \neq u \\ v\in V(G)}}\vert x_v\vert \geq 1+ \vert \sum_{\substack{v \neq u \\ v\in V(G)}}  x_v \vert=2.
				$$
				Now, define $x^{*} = \frac{x}{\Vert x \Vert _1}$. Then $\Vert x^*\Vert _1=1$. As $\sum\limits_{v \in V(G)}x_v=0$, so $ \sum\limits_{v \in V(G)}x^{*}_v=0$. Therefore,
				\begin{align*}
					b(G)&=\min\biggl\{\sum_{uv \in E(G)}\vert y_u - y_v \vert: \sum_{v \in V(G)} y_v =0 ~\mbox{{and}~} \Vert y \Vert_1=1\biggl\}\\
					&\leq \sum\limits_{uv \in E(G)}\vert x^*_u - x^*_v \vert \\
					&=\frac{1}{\Vert x \Vert_1}\sum_{uv \in E(G)} \vert x_u - x_v \vert\\
					&\leq \frac{1}{2}\sum_{uv \in E(G)}\vert x_u - x_v \vert \\
					&\leq \frac{m}{2}\gamma(G).
				\end{align*} 
			\end{proof}
			The following lemma will be useful in the proof of the next theorem. 
			\begin{lemma}\label{norm-2-bound}
				Let $G$ be a connected graph with $n \geq 2$ vertices. Let $x$ be an $l_\infty$-Fiedler vector of $G$. Then $\Vert x \Vert^2 _{2} \geq \frac{n}{n-1}$. Equality holds if and only if $G$ is isomorphic to the complete graph $K_n$.
			\end{lemma}
			\begin{proof}
				Let $u \in V(G)$ such that $x_u=1$. By applying the Cauchy-Schwarz inequality, we obtain 
				$$
				(\sum_{v \neq u}x_v)^2 \leq (\sum_{v \neq u}1^2)(\sum_{v \neq u}x_v^2).
				$$
				Therefore, we have $\sum\limits_{v \neq u}x_v^2 \geq \frac{1}{n-1}$. Hence 
				$\Vert x \Vert^2 _{2} \geq \frac{n}{n-1}$. The equality holds if and only if $x_v=-\frac{1}{n-1}$ for all $v\neq u$. In this case, $\gamma(G)=\gamma_x(G)=\frac{n}{n-1}$. Then, by Theorem \ref{main theorem}, we have $\max\limits_{v \in V(G)}\tr(v)=n-1$. Since $\tr(v) \geq n-1$ for every $v \in V(G)$, we have $\tr(v) = n-1$ for all $v \in V(G)$. Hence, $G$ is isomorphic to $K_n$.
			\end{proof}
			
			Using the bound on the norm of $l_\infty$-Fiedler vector derived above, we now establish a relation between $\gamma (G)$ and the algebraic connectivity $a(G)$.
			\begin{theorem}\label{a(G)-bound}
				Let $G$ be a connected graph with $n\geq 3$ vertices and $m$ edges. Then $a(G) < \frac{m(n-1)}{n}\gamma(G)^2$.
			\end{theorem}
			\begin{proof}
				If $G=K_n$, we have $a(K_n)=n$ \cite{alg-con}. We know $\gamma(K_n)=\frac{n}{n-1}$. Hence,  for $n \geq 3$, we have 
				$a(K_n)=n<\frac{n^2}{2}=\frac{m(n-1)}{n}\gamma(K_n)^2$.
				
				Let $G \neq K_n$. Let $x$ be an $l_\infty$-Fiedler vector of $G$. Observe that, if $vw\in E(G)$, then $\vert x_v - x_w\vert \leq \gamma(G) $. Now, by the variational characterization of  the second smallest eigenvalue\cite{matrix-analysis}, we have
				\begin{align*}
					a(G) &\leq \frac{x^T L(G) x}{x^Tx}\\&= \frac {\sum\limits_{vw \in E(G)} (x_v-x_w)^2}{\vert \vert x \vert \vert^2 _{2}}\\ &< \frac{m(n-1)}{n}\gamma(G)^2,
				\end{align*}
				
				where the last inequality follows from Lemma \ref{norm-2-bound}.
			\end{proof}
			Let $S \subseteq V(G)$. The volume of $S$, denoted by $\vol(S)$, is the sum of the degrees of all vertices in $S$, that is 
			$\vol(S)=\sum\limits_{u \in S} \deg(u).$  For a proper subset $S \subsetneq  V(G)$, the boundary of $S$, denoted by $ \partial S$, is the set of all edges with exactly one endpoint in $S$. Formally 
			$$\partial S = \biggl\{uv \in E(G) : u \in S ~\mbox{and}~v \in S^c\biggl\}.$$ 
			The following definitions and results are from \cite{chung-book-1997}. For a subset $S \subsetneq V(G)$, $h_G(S)$ is defined as:
			$$
			h_G(S)=\frac{\vert \partial S \vert}{\min(\vol(S),\vol(S^c))},
			$$
			where $S^c = V(G) \setminus S.$
			The Cheeger constant $h_G$ of a graph $G$ is given by 
			$$
			h_G= \min\limits_{S \subsetneq V(G)} h_G(S).
			$$
			Let $G$ be a graph on $n$ vertices. Then the normalized Laplacian of $G$, denoted by $\mathcal{L}(G)$, is given by
			
			$$
			\mathcal{L}_{uv}(G) =
			\begin{cases}
				1 & \text{if } u = v \text{ and } \deg(v) \neq 0, \\
				-\dfrac{1}{\sqrt{\deg(u) \deg(v)}} & \text{if u and v are adjacent} , \\
				0 & \text{otherwise}.
			\end{cases}
			$$
			In \cite{chung-book-1997}, the author defined $\mathcal{L}(G)$ to be the Laplacian matrix associated with the graph $G$. If $G$ is connected graph, then we have $\mathcal{L}(G)=\Delta(G)^{\frac{-1}{2}}L(G)\Delta(G)^{\frac{-1}{2}}$ . The spectrum of $\mathcal{L}(G)$ is denoted by $\{ \mu_1 (G),\mu_2 (G),\dots,\mu_n (G)\}$, where the eigenvalues are arranged in non-increasing order: $\mu_1(G) \geq \mu_2(G)\geq \dots \geq \mu_n(G)$.
			The  Cheeger inequality is given by 
			$$
			2h_G \geq \mu_{n-1}(G) > \frac{h_G^2}{2}.
			$$
			\begin{theorem}
				Let $G$ be a connected $k$-regular graph on $n\geq 3$ vertices. Then $h_G < \sqrt{n-1} \gamma(G)$.
			\end{theorem}
			\begin{proof}
				Let $G$ be $k$-regular graph. Then $\mathcal{L}(G)=\frac{1}{k}L(G)$ and  $\mu_{n-1}(G)=\frac{a(G)}{k}$. Since $\vert E(G) \vert =\frac{kn}{2}$, from Theorem \ref{a(G)-bound}, we have $\mu_{n-1}(G) < \frac{n-1}{2} \gamma(G)^2$. The result follows from the Cheeger inequality.
			\end{proof}
			
			
			Note that if $H$ is a graph obtained by adding an edge to the graph $G$, then $\gamma (H) \geq \gamma (G)$. To see this, let $x \in \mathcal{F}$. Then $\gamma_x(H) \geq \gamma_x(G)$. Taking the minimum over all $x \in \mathcal{F}$ on both sides yields the desired inequality.
			
			Next,  we establish bounds on the quantity $\gamma(G)$ in terms of the number of vertices and provide a complete characterization of the extremal graphs.
			\begin{theorem}\label{global extremizer}
				Let $G$ be a connected graph on $n$ vertices. Then
				$$\frac{2}{n-1} \leq \gamma(G) \leq \frac{n}{n-1}.$$
				
				Furthermore,  equality holds on the left-hand side if and only if  $G$ is isomorphic to $P_n$, and equality holds on the right-hand side if and only if  $G$ is isomorphic to $K_n$.  
			\end{theorem}
			\begin{proof}
				%For any vertex $ u \in V(K_n)$, the transmission is $n-1$. 
				If $H$ is a spanning subgraph of a graph $G$, then $\gamma (H) \leq \gamma (G)$. Thus $$\gamma(G) \leq \gamma(K_n) = \frac{n}{n-1}.$$
				
				The upper bound is attained if and only if $\max\limits_{u\in V(G)} \tr(u)=n-1$. Since $\tr(u) \geq n-1$ for every $u \in V(G)$, this implies that $\tr(u) = n-1$ for all $u \in V(G)$. Hence, $G \cong K_n$.\\
				Note that for any vertex $u \in V(G)$,
				$$
				\tr(u) \leq \frac{n(n-1)}{2}.
				$$
				Thus, $$\max\limits_{u \in V(G)} \tr(u) \leq \frac{n(n-1)}{2},$$ which implies  $$\gamma(G) \geq \frac{2}{n-1}.$$
				For equality to hold, the diameter of the graph $G$ should be $n-1$. Hence, $P_n$ is the only graph that attains the lower bound.
			\end{proof}
			
			\begin{theorem}
				Let $T$ be a tree on $n(\geq 3)$ vertices.
				Then 
				$$\frac{2}{n-1} \leq \gamma(T) \leq \frac{n}{2n-3}.$$
				
				Furthermore,  equality holds on the left-hand side if and only if  $G$ is isomorphic to $P_n$, and equality holds on the right-hand side if and only if  $G$ is isomorphic to $S_n$.
			\end{theorem}
			\begin{proof}
				The left-hand side of the inequality follows from Theorem \ref{global extremizer}. By Corollary \ref{star graph}, we have $\mathcal{D}_M(S_n)=2n-3$. We now show that for any tree on $n$ vertices, $\mathcal{D}_M(T) \geq 2n-3$. 
				
				Let $u \in V(T)$ be such that $\tr(u)=\mathcal{D}_M(T)$. Then, by Theorem \ref{pendant}, the vertex $u$ should be a pendant vertex. So, the vertex $u$ has exactly one neighbor, and the other vertices of $T$ are at a distance of at least $2$ from $u$. Hence
				$$
				\mathcal{D}_M(T)=\tr(u)\geq 2n-3.
				$$
				This gives the upper bound: $$\gamma (T) \leq \frac{n}{2n-3}.$$
				To get equality in the above expression, we must have the diameter of the tree $T$ equal to $2$. It is easy to verify that the star graph $S_n$ is the unique tree with diameter $2$. Hence, the star graph is the unique extermal graph.
			\end{proof}
			\section{Cartesian product}\label{sec-product}
			
			
			In this section, we derive a formula for $\gamma(G)$ when $G$ is the Cartesian product of some finite families of graphs, expressing it in terms of the corresponding values for their constituent graphs. As a by-product of this formula, we obtain explicit values of $\gamma(G)$ for several well-known families of graphs.
			
			The \textit{Cartesian product} of two graphs $G$ and $H$, denoted by $G\times H$, is the graph with the vertex set $V(G) \times V(H)$, and edges defined as follows: $(u_1,v_1) \sim (u_2,v_2)$ if either $u_1 \sim u_2$ in $G$ and $v_1=v_2$, or $u_1=u_2$ and $v_1 \sim v_2$ in $H$ \cite{hammack-product-graph}.
			
			
			\begin{theorem}{\label{graph products}}
				
				Let $G$ and $H$ be two connected graphs. Then 
				$$
				\gamma(G \times H)=\frac{1}{\frac{1}{\gamma(G)}+\frac{1}{\gamma(H)}}.
				$$
			\end{theorem}
			\begin{proof}
				Let $G$ and $H$ be graphs of order $n$ and $m$,  respectively. Let $\mathcal{D}_M(G)=t_1$, $\mathcal{D}_M(H)=t_2$. Then $\gamma(G)=n/t_1$ and $\gamma(H)=m/t_2$. It is known that \cite{hammack-product-graph}, for any two vertices $(u_1,v_1)$ and $(u_2,v_2)$ in $G\times H$, 
				
				$$
				d_{G\times H}\{(u_1,v_1),(u_2,v_2)\}=d_G(u_1,u_2)+d_H(v_1,v_2).
				$$
				As a consequence, it is easy to see that the transmission of a vertex $(u,v)$ in $G\times H$ is
				\begin{align*}
					\tr_{G\times H}(u,v)&=m\tr_G(u)+n\tr_H(v).
				\end{align*}
				
				To maximize this over all $(u,v)\in V(G\times H)$, we pick $u \in V(G)$ with transmission $t_1$ and $v \in V(H)$ with transmission $t_2$. Hence,  $\mathcal{D}_M(G \times H)=mt_1+nt_2$. Therefore, by Theorem \ref{main theorem},
				$$
				\gamma(G\times H)=\frac{mn}{mt_1+nt_2}=\frac{1}{\frac{t_1}{n}+\frac{t_2}{m}}=\frac{1}{\frac{1}{\gamma(G)}+\frac{1}{\gamma(H)}}.
				$$
			\end{proof}
			
			
			
			Similarly, the Cartesian product of the graphs $G_1,G_2,\dots,G_k$, denoted by $G_1 \times G_2 \times \dots \times G_k$, is the graph with the vertex set $V(G_1) \times V(G_2) \times \dots \times V(G_k)$, and the edge set defined as follows: $(u_1,u_2,\dots,u_k) \sim (v_1,v_2,\dots,v_k)$ if there exists exactly one index $i \in \{1,2,\dots,k\}$ such that $u_i \sim v_i$ in $G_i$, and $u_j=v_j$ for all $j \neq i$.
			
			
			Building on the proof of the above theorem, we can extend the argument to establish the following result for a finite Cartesian product of graphs.
			
			\begin{theorem}
				Let $G_1, G_2,\dots, G_k$ be a collection of connected graphs. Then
				$$
				\gamma(G_1 \times G_2 \times \dots \times G_k)=\frac{1}{\frac{1}{\gamma(G_1)}+\frac{1}{\gamma(G_2)}+\dots +\frac{1}{\gamma(G_k)}}.
				$$
			\end{theorem}
			
			
			
			Using the above theorem, we can compute the explicit value of $\gamma (G)$ for several well-known families of graphs. We present a few of them here.  
			
			\begin{enumerate}
				\item [(a)] The $t$-dimensional hypercube $Q_t$ is the Cartesian product of $t$-copies of the complete graph $K_2$. Since $\gamma(K_2)=2$, we have $\gamma(Q_t)=\frac{2}{t}$. 
				\item [(b)] The Hamming graph $H(t,s)$ is the Cartesian product of $t$-copies of the complete graph $K_s$. Hence $\gamma(H(t,s))=\frac{s}{t(s-1)}.$
				\item [(c)] The grid graph $G_{l,m,n}$ is the Cartesian product of paths $P_l$, $P_m$ and $P_n$.
				Now, by Theorem \ref{global extremizer}, we obtain $\gamma(G_{l,m,n})=\frac{2}{l+m+n-3}$.
				
				\item[(d)]  The torus grid graph $T_{m,n}$ is the Cartesian product of two        cycle graphs $C_m$ and $C_n$. By Corollary \ref{basic examples}, we have   $\gamma(C_n)=\frac{n}{\lfloor \frac{n^2}{4}\rfloor}$. Thus
				$$
				\gamma(T_{m.n})=\frac{mn}{m\lfloor \frac{n^2}{4} \rfloor +n\lfloor \frac{m^2}{4} \rfloor}.
				$$
			\end{enumerate}
			\section*{Conclusion}
			In this work, we studied the parameter $\gamma(G)$ as an $l_\infty$-analogue of the classical algebraic connectivity and demonstrated that $\gamma(G)$ characterizes the connectedness of the underlying graph $G$. We established a combinatorial formula for $\gamma(G)$ based on the maximum row sum of the distance matrix of $G$, which enables efficient computation of $\gamma(G)$ using a breadth-first search algorithm.  We  completely characterized all $\ell_\infty$-Fiedler vectors.  Additionally, we explored the relationship between $\gamma(G)$ and several other graph parameters. Furthermore, we derived a formula for $\gamma(G)$ when $G$ is the Cartesian product of finitely many graphs, and calculated $\gamma (G)$ for some well-known graphs.
			
			Graph smoothing with respect to the $l_2$-norm is well-studied, while the corresponding problems for the $l_1$-norm and the $l_\infty$-norm were introduced by Andrade and Dahl \cite{enide-geir-2024}. In this paper, we further investigate the graph smoothing problem for the $l_\infty$-norm.  We conclude with the following table summarizing the graph smoothing problems with respect to the $l_2$-norm, $l_1$-norm, and $l_\infty$-norm. The conclusions in the last column are based on the results presented in this work. 
			
			
			\begin{table}[htbp]
				\centering
				\Large % Increase font size
				\resizebox{\textwidth}{!}{%
					\begin{tabular}{|c|c|c|c|}
						\hline
						\textbf{Name} &  $l_2$-graph smoothing  & $l_1$-graph smoothing & $l_\infty$-graph smoothing\\
						\hline
						\textbf{} &   &  & \\
						\textbf{Problem} & $\min\limits_x \sum\limits_{uv \in E(G)} (x_u - x_v)^2$ & $\min\limits_x \sum\limits_{uv \in E(G)} \lvert x_u - x_v \rvert$ & $\min\limits_x \max\limits_{uv \in E(G)} \lvert x_u - x_v \rvert$ \\
						& subject to $\| x \|_2 = 1$, $x \perp e$ & subject to $\| x \|_1 = 1$, $x \perp e$ & subject to $\| x \|_\infty = 1$, $x \perp e$ \\
						\textbf{} &   &  & \\
						\hline
						\textbf{Solution} & $a(G)$ & $b(G)$ & $\gamma (G)$ \\
						\hline
						\textbf{Encodes} & Yes & Yes & Yes \\
						\textbf{the connectedness} &  &  &  \\
						\hline
						\textbf{Relevance} & partition & Sparsest cut &  max transmission \\
						\hline
						\textbf{Strategy} & spectral & combinatorial & combinatorial \\
						&  & (NP-hard) & (Polynomial-time) \\
						\hline
					\end{tabular}%
				}
				\caption{\text{Comparison of graph smoothing problems under different norms.}}
				\label{tab:graph_smoothing}
			\end{table}
			
			
			\section*{Acknowledgments}
			M. Rajesh Kannan acknowledges financial support from the ANRF-CRG India and  SRC, IIT Hyderabad. Rahul Roy thanks the University Grants Commission (UGC), India, for financial support. The authors thank ChatGPT for helpful discussions on improving the exposition of the manuscript, identifying typographical inconsistencies, and refinements of proofs. The authors have independently verified all mathematical statements and are solely responsible for the correctness of the results presented in this paper.
			
			
			
			
			
			\bibliographystyle{amsplain}
			\bibliography{gamma(X)_ref}
			
			
			\vspace{0.4cm}
			
			\affl{M. Rajesh Kannan}{rajeshkannan@math.iith.ac.in, rajeshkannan1.m@gmail.com}{Department of Mathematics, Indian Institute of Technology Hyderabad, Kandi, Sangareddy 502284, India.}
			
			\affl{Rahul Roy}{ ma23resch11004@iith.ac.in}{Department of Mathematics, Indian Institute of Technology Hyderabad, Kandi, Sangareddy 502284, India.}
			
			
			
		\end{document}
		
